% BHU PYQ -- Manually transcribed & error-corrected
\documentclass[11pt,a4paper]{article}

\usepackage[a4paper,top=2.5cm,bottom=2.5cm,left=2.8cm,right=2.8cm,headheight=14pt]{geometry}
\usepackage{amsmath,amssymb}
\usepackage[T1]{fontenc}
\usepackage[utf8]{inputenc}
\usepackage{lmodern}
\usepackage{microtype}
\usepackage{fancyhdr}
\usepackage{enumitem}
\usepackage{hyperref}
\usepackage{lastpage}
\usepackage{parskip}

\hypersetup{colorlinks=true,urlcolor=black,linkcolor=black,
  pdftitle={CHB-604: Physical Chemistry-IV -- BHU PYQ},pdfauthor={Banaras Hindu University}}

\pagestyle{fancy}\fancyhf{}
\renewcommand{\headrulewidth}{0.4pt}\renewcommand{\footrulewidth}{0.4pt}
\fancyhead[L]{\small   \textit{Banaras Hindu University}}
\fancyhead[R]{\small   \textit{CHB-604: Physical Chemistry-IV}}
\fancyfoot[C]{\small Page  \thepage\ of \pageref{LastPage}}


\newlist{parts}{enumerate}{2}
\setlist[parts,1]{label=(\alph*),leftmargin=*,itemsep=5pt,topsep=3pt}
\setlist[parts,2]{label=(\roman*),leftmargin=*,itemsep=3pt,topsep=2pt}

\newcommand{\pts}[1]{\hfill   \textit{[#1]}}


\begin{document}


\begin{center}
  {\large\bfseries BANARAS HINDU UNIVERSITY}\\[2pt]
  {\normalsize B.Sc. (Hons.) Semester VI Examination, 2013-14}\\[6pt]
  \rule{0.8\linewidth}{0.5pt}\\[6pt]
  {\large\bfseries Paper No. CHB-604: Physical Chemistry-IV}\\[4pt]
  {\normalsize Time: 3 Hours\hfill Maximum Marks: 70}
\end{center}
\rule{\linewidth}{0.4pt}
\small



\noindent   \textit{Instructions: Answer five questions in total, including Question~1 which is compulsory.}

\normalsize
\bigskip




\noindent   \textbf{Question 1.} Answer all parts of the following: \pts{5 $ \times$ 2 = 10}

\begin{parts}
  \item Protons of $ \text{CH}_3 \text{I}$ appear at $130\, \text{Hz}$ in a $50\, \text{MHz}$ NMR instrument. At what frequency would these be observed in a $600\, \text{MHz}$ instrument?
  \item Explain the conditions a function should satisfy for becoming a well-behaved wave function in quantum chemistry.
  \item Differentiate between 'state' and 'microstate' in statistical thermodynamics.
  \item Write the expression for the vibrational partition function.
  \item Differentiate between Fermi resonance and coupling in IR spectroscopy.
\end{parts}

\medskip\rule{\linewidth}{0.2pt}




\noindent   \textbf{Question 2.}

\begin{parts}
  \item Write the Schr\"odinger wave equation for a one-dimensional harmonic oscillator. Write general expressions for its energy levels and wave functions. Sketch the lowest two wave functions. \pts{8}
  \item Write the time-dependent Schr\"odinger wave equation. Under which conditions can you transform it into the time-independent equation? \pts{6}
\end{parts}

\medskip\rule{\linewidth}{0.2pt}




\noindent   \textbf{Question 3.}

\begin{parts}
  \item What are the factors that determine the intensity of a spectral line? \pts{4}
  \item How would you classify rotating molecules on the basis of their three principal moments of inertia ($I_A, I_B, I_C$)? \pts{4}
  \item How would you evaluate the two bond lengths in the linear $ \text{OCS}$ molecule using microwave spectral data? \pts{6}
\end{parts}

\medskip\rule{\linewidth}{0.2pt}




\noindent   \textbf{Question 4.}

\begin{parts}
  \item Show that for an anharmonic oscillator, the selection rules are $\Delta v = \pm1, \pm2, \pm3, \dots$. Why is only the transition originating from $v=0$ considered? \pts{6}
  \item State the Franck-Condon principle and illustrate it using potential energy curves for the ground and excited electronic states. \pts{8}
\end{parts}

\medskip\rule{\linewidth}{0.2pt}




\noindent   \textbf{Question 5.}

\begin{parts}
  \item What are the spin-lattice ($T_1$) and spin-spin ($T_2$) relaxation times in NMR spectroscopy? Explain their physical meanings. \pts{6}
  \item Define the relation $S = k_B \ln \Omega$. Using this, derive the translational molecular partition function of an ideal gas. \pts{8}
\end{parts}

\vfill
\rule{\linewidth}{0.4pt}

\begin{center}\small
  \href{https://bhu-exam.vercel.app/}{Click here to visit our website}\\[4pt]
  \href{https://me-aryan.vercel.app/}{Click here to visit our contributor}\\[4pt]
  \href{https://quashan-me.vercel.app/}{Click here to visit our contributor}
\end{center}

\end{document}
